% [原文 p.67]
\chapter{泰希米勒空间}

\ChapterAbstract{光滑曲面 $S$ 的泰希米勒空间，是 $S$ 上双曲度量模去同痕于恒等映射的微分同胚所成的空间。等价地，它在同一等价关系下参数化 $S$ 上的共形（或复）结构。通过和乐表示，泰希米勒空间也可等同于 $\PSL_2(\R)$ 特征簇的一个连通分支。本章我们给出主要定义，并显式写出和乐表示。}

由于泰希米勒空间刻画曲面上的度量，它们在物理的诸多领域中自然扮演重要角色，其中最突出的是二维与三维引力。此外，在研究 class $\mathcal{S}$ 的 $4d$ $\mathcal{N}=2$ 超对称量子场论时，它们占有格外特殊的地位：在这些理论中，它们最自然地解释为光滑曲面上复结构的参数空间。泰希米勒空间在这一语境中的出现将在第 21 章更详细地讨论。

\section{定义}

数学中的一大主题，是把对象按某种等价关系分类。一个简单的例子是把可定向紧曲面按同胚分类。这样的曲面由亏格 $g\in\N$ 与边界分支数 $h$ 刻画。泰希米勒空间通过加入几何来细化这一拓扑分类。更具体地说，它们对带有常曲率度量的曲面进行分类。这里我们对泰希米勒理论只作极为简短的浏览。更多细节，我们推荐 Hubbard [211]、Farb–Margalit [125] 的书，以及 Papadopoulos 主编的系列丛书 [295]。

任何曲面都可赋予常曲率度量（例如来自其万有覆叠，后者或是 $\mathbb{S}^2$、$\R^2$，或是 $\HH^2$），并使边界为测地线。Gauss–Bonnet 定理把曲率常数 $K$ 与拓扑联系起来：欧拉示性数为正时 $K>0$，为零时 $K=0$，为负时 $K<0$。

% [原文 p.68]
我们关心最后一种情形，因此考虑紧、连通、可定向的拓扑曲面 $\bar S_{g,h}$，其亏格为 $g$、有 $h$ 个边界分支，并满足 $\chi(\bar S_{g,h})=2-2g-h<0$（这样的曲面也称为双曲型的）。于是任何曲面 $\bar S_{g,h}$ 都容许一个双曲结构，如下所定义：

\begin{definition}[6.1] 如上拓扑曲面 $\bar S_{g,h}$ 上的一个双曲结构是一个二元组 $(M,f)$，其中 $M$ 是双曲曲面，即带有常负曲率 $K=-1$ 的度量与测地边界的拓扑曲面，而 $f:\bar S_{g,h}\to M$ 是保向同胚。\end{definition}

映射 $f:\bar S_{g,h}\to M$ 称为标记（marking）。它负责给曲面 $\bar S_{g,h}$ 装备一个常负曲率 $-1$ 的黎曼度量。于是泰希米勒空间就是这些二元组在某个等价关系下的商空间：

\begin{definition}[6.2] $\bar S_{g,h}$ 的泰希米勒空间 $\mathcal{T}(\bar S_{g,h})$ 是空间
\[
\mathcal{T}(\bar S_{g,h}):=\{(M,f)\}/\sim,
\]
其中
\begin{itemize}
  \item $(M,f)$ 是 $\bar S_{g,h}$ 上的双曲结构；
  \item $(M,f)\sim(M',f')$ 当且仅当存在等距 $\iota:M\to M'$，使得 $\iota\circ f$ 同伦于 $f'$；该等价关系称为同痕（isotopy）。
\end{itemize}
\end{definition}

令人惊讶的是，除双曲结构之外，曲面上还有若干其他几何结构，例如共形结构与复结构，它们也由泰希米勒空间参数化。这来自著名的

\begin{theorem}[6.3（庞加莱单值化）] 任何单连通黎曼曲面都双全纯于 $\C\PP^1$、$\C$ 或双曲平面 $\HH^2$ 之一。\end{theorem}

利用双倍论证（doubling argument），可以把单值化定理的结果推广到具有非空边界的黎曼曲面：

\begin{theorem}[6.4] 设 $\bar S_{g,h}$ 是具有非空边界的黎曼曲面。则 $\bar S_{g,h}$ 的万有覆叠双全纯于 $\C$ 的一个具有非空内部的单连通闭真子集。\end{theorem}

由此可知，$\bar S_{g,h}$ 是双曲型的当且仅当其万有覆叠双全纯于 $\HH^2$ 的一个闭子集。

% [Tutor Note: 定理 6.4 的“双倍论证”细节：沿全部边界分量把带边曲面 $\bar S_{g,h}$ 双倍，得到闭黎曼曲面 $D$，对 $D$ 用定理 6.3；$D$ 上交换两半的对合 $\sigma$ 的不动点集正是原边界，在万有覆叠上取提升后的不动点分支，即得 $\bar S_{g,h}$ 的万有覆叠。注意定理 6.4 隐含了 $\chi<0$（双曲型）这一前提：若 $\chi=0$，例如 $g=0,h=2$ 的环带，其万有覆叠双全纯于整个 $\C$，并不是 $\C$ 的真子集。原文在本段之前已声明只关心 $\chi<0$ 的情形，故定理应在此前提下理解。]

考虑 $\bar S_{g,h}$ 上的复结构，即一个图册，其坐标卡取值于 $\C$、转移函数全纯。则其万有覆叠 $\tilde S_{g,h}$ 也带有复结构，因为局部地 $\bar S_{g,h}$ 与 $\tilde S_{g,h}$ 同胚。我们通过用等距作用的覆叠变换取商，从 $\tilde S_{g,h}\subset\HH^2$ 恢复出 $\bar S_{g,h}$。因此，对 $\bar S_{g,h}$ 上的任何复结构，我们都可联系一个 $\bar S_{g,h}$ 上的双曲结构；反之，所有双曲结构都如此产生。

% [原文 p.69]
\begin{theorem}[6.5] 泰希米勒空间 $\mathcal{T}(\bar S_{g,h})$ 是复结构模去 $\bar S_{g,h}$ 的同痕于恒等映射的微分同胚所成的空间。\end{theorem}

纯映射类群（pure mapping class group）$\operatorname{PMod}(\bar S_{g,h})$，即 $\bar S_{g,h}$ 的所有保向、逐点固定各边界分支的微分同胚，模去相对边界同痕于恒等映射者所成的群，通过改变标记自然地作用在 $\mathcal{T}(\bar S_{g,h})$ 上；该作用真不连续（properly discontinuous），其商称为黎曼曲面的模空间 $\mathcal{M}(\bar S_{g,h})$，这些黎曼曲面的底层拓扑曲面为 $\bar S_{g,h}$。这里“真不连续”是指：对每个紧子集 $K\subset\mathcal{T}(\bar S_{g,h})$，集合
\[
\{g\in\operatorname{PMod}(\bar S_{g,h}) : gK\cap K\neq\varnothing\}
\]
是有限的。

关于泰希米勒空间的一个著名结果如下：

\begin{theorem}[6.6] 紧连通可定向曲面 $\bar S_{g,h}$ 的泰希米勒空间 $\mathcal{T}(\bar S_{g,h})$ 同胚于实维数 $6g-6+3h$ 的开球。\end{theorem}

证明该定理的一个直接方法，是用 Fenchel–Nielsen 坐标参数化 $\mathcal{T}(\bar S_{g,h})$，这将在第 7 章更仔细地研究；完整的证明可在 [309, Thm. 9.7.4] 中找到。

% [Tutor Note: 定理 6.6 的维数 $6g-6+3h$ 可用第 7 章的 Fenchel–Nielsen 坐标核对：$3g-3+h$ 条裤子曲线各给一个长度与一个扭转，共 $2(3g-3+h)$ 个参数；再加 $h$ 条边界分量的双曲长度，合计 $2(3g-3+h)+h=6g-6+3h$。极限情形：闭曲面 $h=0$ 退回 $6g-6$；裤子（$g=0,h=3$）得 $3$，恰是三条边界长度。注意：若把边界长度固定（只看模空间），维数是 $6g-6+2h$，别与定理 6.6 混用。]

\section{和乐表示}

% [原文 p.69]
泰希米勒空间 $\mathcal{T}(\bar S_{g,h})$ 的另一种代数实现可以借助和乐表示来设想，下面对此加以描述。这使得我们可以纯粹用基本群 $\pi_1(\bar S_{g,h})$ 及其到上半平面 $\HH^2$ 的保向等距群 $\PSL_2(\R)$ 的表示来理解 $\mathcal{T}(\bar S_{g,h})$。

设 $(M,f)$ 是曲面 $\bar S_{g,h}$ 上的双曲结构。固定基点后，同胚 $f:\bar S_{g,h}\to M$ 诱导基本群的同构
\[
f_*:\pi_1(\bar S_{g,h})\to\pi_1(M).
\]
此外，群 $\pi_1(M)$ 作为覆叠变换群以保向等距作用在万有覆叠 $\tilde M$ 上。但 $M$ 是双曲曲面，由庞加莱单值化定理知 $\tilde M$ 双全纯于 $\HH^2$ 的闭子集，而群 $\PSL_2(\R)$ 恰是保向等距群 $\operatorname{Isom}^+(\HH^2)$。因此我们推得，群 $\pi_1(M)$ 作为离散、余紧的子群嵌入群 $\PSL_2(\R)$；拓扑空间上的群作用若其相伴随的商空间是紧空间，则称为余紧的（cocompact）。

综上，任何双曲结构 $(M,f)$ 都定义一个单射同态 $\rho:\pi_1(\bar S_{g,h})\to\PSL_2(\R)$，其像离散且余紧。此外可以证明，由等价双曲结构诱导的这类表示恰在 $\PSL_2(\R)$ 的元素下共轭，反之亦然。于是我们实现了：
\[
\mathcal{T}(\bar S_{g,h})\subset\Hom\bigl(\pi_1(\bar S_{g,h}),\PSL_2(\R)\bigr)/\PSL_2(\R).
\]

% [原文 p.70]
当 $h=0$，即 $\bar S_{g,h}=S_g$ 为闭曲面时，下一定理说明实际上可以说得更多（见 [167]、[198]）：

\begin{theorem}[6.7] 闭曲面 $S_g$ 的泰希米勒空间 $\mathcal{T}(S_g)$ 是表示空间 $\Hom\bigl(\pi_1(S_g),\PSL_2(\R)\bigr)/\PSL_2(\R)$ 的一个连通分支。\end{theorem}

\begin{remark}[6.8] 事实上，当 $\bar S_{g,h}=S_g$ 为闭曲面时，$\mathcal{T}(S_g)$ 是空间 $\Hom\bigl(\pi_1(S_g),\PSL_2(\R)\bigr)/\PSL_2(\R)$ 的两个完全由离散忠实表示组成的连通分支之一。另一个这样的分支是 $\mathcal{T}(S_g^{\mathrm{op}})$，其中 $S_g^{\mathrm{op}}$ 表示反向定向的曲面。从此观点看，作为离散忠实表示的集合，泰希米勒空间最早出现在 Fricke 与 Klein [140] 的工作中，他们用 Fuchsian 群来定义它。这一进路将促使我们在第 9.1 节引入特征簇。\end{remark}

对泰希米勒空间的研究，因坐标的存在而变得便利，我们下面就来描述这些坐标。

% [原文 p.71]
\chapter{泰希米勒空间上的坐标}

\ChapterAbstract{泰希米勒空间可通过与裤子分解相联系的 Fenchel–Nielsen 长度参数与扭转参数来坐标化。对纤毛曲面，也可以改用 Thurston 的剪切坐标，其变换规则结构格外良好，并自然地引出簇代数结构。}

\section{Fenchel–Nielsen 坐标}

对一条裤子 $P$ 与正实数 $a,b,c\in\R_{>0}$，$P$ 上存在唯一的双曲结构，使得 $P$ 的边界分支是分别具有双曲长度 $a$、$b$、$c$ 的测地线（见第 4.4 节）。给定两条双曲裤子 $P$ 与 $P'$，二者都有一个双曲长度为 $a\in\R_{>0}$ 的边界分支（分别记为 $C$ 与 $C'$），可以通过等同 $C$ 与 $C'$ 把 $P$ 与 $P'$ 粘起来。使得两个原曲面上的双曲度量光滑地粘合起来的做法，构成一个单参数族的互不等价方式。这个实参数称为所得曲面中对应于简单闭曲线 $C=C'$ 的扭转（twist）。也可以把同一条双曲裤子的不同边界分支粘起来，如图 7.1 所示。

设 $S_g$ 是亏格为 $g$ 的闭、连通、可定向拓扑曲面。确定一个裤子分解恰好需要 $3g-3$ 条简单闭曲线 $\gamma_1,\dots,\gamma_{3g-3}$，而该分解包含 $2g-2$ 条裤子。给定双曲结构 $\kappa\in\mathcal{T}(S_g)$，这些简单闭曲线 $\gamma_i$ 各自被赋予其自由同伦类中唯一测地代表的双曲长度 $l_i$。此外，选取一族简单闭曲线 $\nu_1,\dots,\nu_r$（称为标记），它们把每条裤子分解成两个六边形。注意数 $r$ 不是先验固定的，而依赖于 $\nu_1,\dots,\nu_r$ 的选择。这一额外选择为每个 $\gamma_i$ 提供了沿其的扭转参数 $\theta_i\in\R$（对 $\kappa\in\mathcal{T}(S_g)$），因此 Fenchel–Nielsen 坐标由映射
\begin{equation}
\begin{aligned}
FN:\mathcal{T}(S_g)&\longrightarrow \R^{3g-3}_{>0}\times\R^{3g-3},\\
\kappa&\longmapsto (l_1,\dots,l_{3g-3},\theta_1,\dots,\theta_{3g-3}).
\end{aligned}
\tag{7.1}
\end{equation}
定义。

% [Tutor Note: 原文记号不一致：本页把标记曲线族记作 $\nu_1,\dots,\nu_r$，而 p.72 及定理 7.2 把同一对象写作 $\nu_1,\dots,\nu_k$。这不影响 (7.1)：$FN$ 的分量数恒为 $6g-6$，与标记曲线条数无关——标记只用来固定每条裤子如何被剖成两个六边形，从而固定扭转参数的零点；其条数 $r$（或 $k$）随裤子分解与剖分的不同而变，故定理 7.2 说 $FN$ 只依赖裤子曲线与标记的“选择”。照原文译出，读者按“标记曲线族”统一理解即可。]

% [原文 p.72]
\Fig[0.85\linewidth]{figures/fig_7_1.png}{7.1}{粘合一条双曲裤子的两个边界分支。}

亏格 $g=3$ 的闭可定向光滑曲面的裤子分解的一个例子见图 7.2 左图。

\Fig[0.95\linewidth]{figures/fig_7_2.png}{7.2}{双曲长度记为 $l_i$ 与 $k_i$，扭转记为 $\phi_i$ 与 $\psi_i$。}

不失一般性，我们取 $i=1$，并定义沿 $\gamma_1$ 的扭转参数 $\theta_1$（这里我们遵循 [261, Sect. 7.3.1] 的讨论）。取一个双曲结构 $\kappa\in\mathcal{T}(S_g)$，并假设所有 $\gamma_i$ 关于 $\kappa$ 都是测地线。环路 $\gamma_1$ 分隔两条（未必不同的）测地裤子（见图 7.3）。我们记 $S'$ 为 $S_g$ 中由这两条沿 $\gamma_1$ 粘合的裤子之并所成的子曲面。设 $\lambda$ 与 $\lambda'$ 为族 $(\nu_1,\dots,\nu_k)$ 中与 $\gamma_1$ 相交的曲线限制到 $S'$ 所得的两条曲线。不失一般性，我们取 $\lambda$。令 $p:=\gamma_1\cap\nu_1$。我们固定一个提升 $\tilde p\in\tilde S=\HH^2$，并提升所有与 $p$ 关联的曲线：测地环路 $\gamma_1$ 提升为测地线 $\tilde\gamma_1$，而（非测地的）路径 $\lambda$ 提升为一条曲线 $\tilde\lambda$，它连接闭测地线 $\gamma_2$ 与 $\gamma_3$ 的两个提升 $\tilde\gamma_2$ 与 $\tilde\gamma_3$（见图 7.4）。

% [原文 p.73]
\Fig[0.55\linewidth]{figures/fig_7_3.png}{7.3}{子曲面 $S'$，它是两条沿 $\gamma_1$ 粘合的双曲裤子之并。}

\Fig[0.42\linewidth]{figures/fig_7_4.png}{7.4}{子曲面 $S'$ 到万有覆叠的一个提升。}

现在我们取唯一的两条正交测地线 $\tilde\eta_2$ 与 $\tilde\eta_3$，它们分别连接 $\tilde\gamma_1$ 与 $\tilde\gamma_2$、$\tilde\gamma_3$。

\begin{definition}[7.1] $\tilde\eta_2$ 与 $\tilde\eta_3$ 之间的剪切（shear）$s_1\in\R$，是 $\tilde\gamma_1$ 上介于交点 $\tilde\gamma_1\cap\tilde\eta_2$ 与 $\tilde\gamma_1\cap\tilde\eta_3$ 之间的测地线段的带号长度。若沿 $\gamma_1$ 行走的观察者先在其左侧遇到另一条测地线（如图 7.4），则 $s_1$ 的符号为正。\end{definition}

最后，我们把沿 $\gamma_1$ 的扭转参数 $\theta_1$ 定义为 $\theta_1:=\frac{2\pi s_1}{l_1}$。

下一定理表明 (7.1) 中的映射 $FN$ 良定义，并给出 $\mathcal{T}(S_g)$ 的一个坐标化：

\begin{theorem}[7.2（Fenchel–Nielsen 坐标，{{{{{{[261, Thm. 7.3.1]}}}}}}）] 映射 $FN$ 良定义、双射，且只依赖于裤子曲线 $\gamma_1,\dots,\gamma_{3g-3}$ 与标记 $\nu_1,\dots,\nu_k$ 的选择。\end{theorem}

% [Tutor Note: 定义 7.1 的符号约定依赖 $\gamma_1$ 的定向：“沿 $\gamma_1$ 行走”应理解为沿 $\tilde\gamma_1$ 的一个取定方向行走；$\gamma_1$ 定向反转会使 $s_1$ 变号，从而 $\theta_1=2\pi s_1/l_1$ 变号。这与 Fenchel–Nielsen 坐标的通常约定一致：$\theta_1$ 是模 $2\pi$ 的扭转角，改变标记（或改变 $\nu_1$ 的选取）只平移 $\theta_1$。核对方法：在图 7.4 中 $s_1$ 是两条正交测地线 $\tilde\eta_2,\tilde\eta_3$ 在 $\tilde\gamma_1$ 上两交点之间的带号距离，而不是弧长参数之差。]

% [原文 p.74]
同一曲面一般有许多不同的带标记裤子分解，每一个都给出曲面 $S_g$ 的泰希米勒空间 $\mathcal{T}(S_g)$ 上不同的一组 Fenchel–Nielsen 坐标，它们可解释为 $\mathcal{T}(S_g)$ 上不同的整体坐标卡。人们可以计算从一个坐标卡到另一个坐标卡的转移函数，但这些转移函数不具有好的代数性质。

对于由亏格 $g\ge0$ 的可定向闭曲面挖去 $h>0$ 个开圆盘所得的有限型紧曲面 $\bar S_{g,h}$（其中 $g$ 与 $h$ 满足 $\chi(\bar S_{g,h})=2-2g-h<0$），除 $k$ 条边界环路外，还需要 $3g-3+h$ 条简单闭曲线来定义一个裤子分解，而该分解涉及 $2g-2+h$ 条裤子。给定双曲结构 $\kappa\in\mathcal{T}(\bar S_{g,h})$，$\bar S_{g,h}$ 的边界分支是双曲测地线，完全由它们的长度描述，因此 Fenchel–Nielsen 坐标给出映射
\begin{equation}
FN:\mathcal{T}(\bar S_{g,h})\longrightarrow \R^{3g-3+2h}_{>0}\times\R^{3g-3+h}.
\tag{7.2}
\end{equation}
对应于 $g=2$、$h=1$ 的一个例子见图 7.2 右图。

我们可以扩大映射 $FN$ 的定义域，允许 $\bar S_{g,h}$ 的边界分支退化为尖点（cusp）。为使陈述精确，我们不考虑紧曲面 $\bar S_{g,h}$，而取具有 $h$ 个内部穿孔且无边界的纤毛曲面 $S_{g,h}$。我们可以给这样的曲面装备一个（可能不完备的）有限面积凸双曲度量。我们把所有这类度量模去同痕所成的空间记为 $\mathcal{T}(S_{g,h})$。于是对每个穿孔，$\mathcal{T}(S_{g,h})$ 上都有良定义的长度函数，它是环绕该穿孔的所有简单环路长度的下确界。这些长度函数取值于闭半射线 $\R_{\ge0}$。注意一个穿孔是尖点当且仅当其长度等于零。于是 Fenchel–Nielsen 映射也能参数化空间 $\mathcal{T}(S_{g,h})$：
\[
FN:\mathcal{T}(S_{g,h})\longrightarrow (\R_{>0}\times\R)^{3g-3+h}\times\R^{h}_{\ge0}.
\]

% [Tutor Note: 本段原文写“除 $k$ 条边界环路外”，但此处讨论的 $\bar S_{g,h}$ 只有 $h$ 条边界分支，且本节并未定义 $k$（$k$ 仅是标记曲线族的条数记号）。按语境应为“除 $h$ 条边界环路外”，即 $3g-3+h=(3g-3)+h$ 中的 $h$ 恰是边界环路。原文如此印，照译并在此指出。]

注意，若取定 $\mathcal{T}(S_{g,h})$ 中的一个双曲结构，则对每个正长度 $l>0$ 的穿孔，该穿孔的一个同胚于开圆环的邻域可以通过添加一条长度为 $l$ 的闭测地线来紧化，这对应于该圆环的度量完备化。若所有穿孔都有正长度，则这给出双曲纤毛曲面 $S_{g,h}$ 的一个紧化，即第 4.2 节意义下的紧双曲曲面 $\bar S_{g,h}$，它由 (7.2) 参数化。

% [Tutor Note: 式 (7.2) 的像空间维数可自校验：$(3g-3+2h)+(3g-3+h)=6g-6+3h$，与定理 6.6 的维数完全一致。其中 $3g-3+2h$ 由 $3g-3+h$ 条内部裤子曲线的长度与 $h$ 条边界分量的长度组成，而 $3g-3+h$ 是扭转数；这与 p.71 的闭曲面情形 $(3g-3)+(3g-3)=6g-6$ 相容。]

\section{剪切坐标}

本节我们将在纤毛曲面 $S:=S_{g,p}$ 的泰希米勒空间的一个略有不同的版本上引入另一组坐标。

回忆纤毛曲面 $S$ 可装备一个带标记的、可能不完备的、凸的、有限面积的双曲结构，可能带边界，使得纤毛之间的边界段是无穷测地线，且纤毛是无穷远点。此外，每个内部穿孔都联系一个长度（见定义 4.9）。

% [原文 p.75]
每个正长度穿孔的度量完备化给曲面添加一条闭边界测地线。注意这些边界测地线各有两种不同的定向。

\begin{definition}[7.3] 我们记 $\mathcal{T}^x(S)$ 为 $S$ 上如上所有带标记双曲结构模去同痕所成的空间，并附带对每个正长度穿孔的度量完备化所对应边界测地线之定向的一个选择。我们称空间 $\mathcal{T}^x(S)$ 为纤毛曲面 $S$ 的带标架泰希米勒空间（framed Teichmüller space）。\end{definition}

\begin{remark}[7.4] 注意“带洞曲面的泰希米勒空间”这一概念最早在 [128] 中引入（并在 [129]、[131] 中进一步研究），目的是描述具有紧测地边界分支的曲面上所有双曲结构所成的空间，并附带边界测地线定向的选择。我们的语境在某种意义上略微更一般：我们考虑带穿孔与纤毛的曲面。这一约定在稍后第 10.2 节研究带标架基本群表示与正性概念时会证明是有用的。\end{remark}

本节我们将定义的 $\mathcal{T}^x(S)$ 上的整体坐标系，与 $S$ 的三角剖分相联系（见第 4.3 节）。给定 $S$ 的一个三角剖分 $\Delta$ 与一个双曲结构 $\kappa\in\mathcal{T}^x(S)$，可以给 $\Delta$ 的每条内部边联系一个实数 $s_\gamma$；这定义了一个称为剪切坐标的映射
\[
\begin{aligned}
Sh:\mathcal{T}^x(S)&\to\R^{E_i(\Delta)}\\
\kappa&\mapsto (s_\gamma)_{\gamma\in E_i(\Delta)},
\end{aligned}
\]
其中 $E_i(\Delta)$ 是 $\Delta$ 的内部边集。

给定双曲结构 $\kappa\in\mathcal{T}^x(S)$，每个理想三角剖分都可变成测地理想三角剖分，即边为测地线的三角剖分：这些测地边或连接无穷远点（尖点与纤毛），或对正长度穿孔按 $\kappa$ 规定的方向绕边界测地线盘旋。理想三角剖分的每个面都是一个理想测地双曲三角形。

$\Delta$ 提升到 $S$ 的万有覆叠 $\tilde S\subseteq\HH^2$，得到凸区域 $\tilde S$ 的一个无穷 $\pi_1(S)$-不变理想测地三角剖分 $\tilde\Delta$。我们考虑 $\tilde\Delta$ 中一个顶点为 $x_1,x_2,x_3,x_4\in\partial_\infty\tilde S$ 的理想四边形，它由连接 $x_1$ 与 $x_3$ 的边 $\gamma$ 三角剖分（见图 7.5）。

现在我们取唯一的两条正交测地线 $\tilde\eta_2$ 与 $\tilde\eta_4$，它们分别连接 $\tilde\gamma$ 与 $x_2$、$x_4$。我们以这种方式把按定义 7.1 定义的 $\tilde\eta_2$ 与 $\tilde\eta_4$ 的剪切 $s_\gamma$ 联系到边 $\gamma$。由于该构造是 $\pi_1(S)$-不变的，$\tilde\Delta$ 上的坐标下降为 $\Delta$ 上的坐标。

\begin{theorem}[7.5（剪切坐标，{{{{{{[261, Prop. 7.4.9]}}}}}}）] 映射 $Sh$ 良定义、双射，且只依赖于三角剖分 $\Delta$ 的选择。\end{theorem}

给定双曲结构 $\kappa\in\mathcal{T}^x(S)$，用剪切坐标容易判断 $S$ 上的一个穿孔是尖点还是容许一条边界测地线，在后一情形下还容易判断它是如何定向的。

% [原文 p.76]
\Fig[0.9\linewidth]{figures/fig_7_5.png}{7.5}{三角剖分四边形在 $\HH$ 中的标准形。}

\begin{proposition}[7.6（{{{{{{[261, Sects. 7.4.4-7.4.5]}}}}}}）] 设 $\gamma_1,\dots,\gamma_k$ 是 $\Delta$ 中所有与 $S$ 的穿孔 $p$ 关联的边，$\kappa\in\mathcal{T}^x(S)$，并设 $c_i$ 是联系到 $\gamma_i$ 的剪切坐标（$i\in\{1,\dots,k\}$）。数 $C=\sum_{i=1}^k c_i$ 为零当且仅当 $p$ 是尖点。$C>0$ 当且仅当 $p$ 的长度为 $C$，且相应的边界测地线的定向与 $S$ 的定向一致。$C<0$ 当且仅当 $p$ 的长度为 $-C$，且相应的边界测地线的定向与 $S$ 的定向相反。\end{proposition}

剪切坐标可以解释为交比。由于点 $x_1,\dots,x_4\in\partial_\infty\tilde S\subseteq\partial_\infty\HH^2\cong\R\PP^1$，我们可以计算交比：
\begin{equation}
z_\gamma:=z:=[x_1,x_2,x_3,x_4]=-\frac{(x_1-x_2)(x_3-x_4)}{(x_2-x_3)(x_4-x_1)}.
\tag{7.3}
\end{equation}
直接计算表明 $z_\gamma=\ee^{s_\gamma}$。

数 $z$ 还可以如下解释：由于把 $(x_1,x_2,x_3)$ 映到 $(\infty,-1,0)$ 的 $\PSL_2(\R)$ 元素是唯一的，$z$ 由循环有序四元组 $(x_1,x_2,x_3,x_4)$ 以及首元素 $x_1$ 的选择唯一确定。对循环集 $(x_1,x_2,x_3,x_4)$ 而言，满足“四元组的前三个元素是某个三角形的顶点”这一三角剖分相容性条件的另一可能首元素选择是 $x_3$。若代替四元组 $(x_1,x_2,x_3,x_4)$ 而取 $(x_3,x_4,x_1,x_2)$，则被映到 $\infty,-1,0$ 的是顶点为 $(x_3,x_4,x_1)$ 的三角形，而第二个三角形的顶点在 $x_1,x_2,x_3$。一个非常令人愉快的事实（它可以通过显式计算验证，也是式 (7.3) 的直接推论）是：两种情形下的剪切参数实际上相同。因此我们定义的剪切参数是 $\PSL_2(\R)$-不变的，且只依赖于四边形的顶点 $x_1,x_2,x_3,x_4$ 以及它被如何三角剖分。

% [Tutor Note: 式 (7.3) 的归一化与图 7.5 右图吻合：用 $\PSL_2(\R)$ 把 $(x_1,x_2,x_3)$ 送到 $(\infty,-1,0)$ 后，把 $x_1=\infty$ 代入 (7.3) 得 $[\infty,-1,0,z]=\frac{x_3-x_4}{x_2-x_3}=\frac{0-z}{-1-0}=z$，即 $z$ 正是 $x_4$ 的像——与图 7.5 右图把 $x_4$ 记为 $z$ 一致。这也解释了负号的作用：去掉负号会得到 $-z$，与下一句 $z_\gamma=\ee^{s_\gamma}>0$ 矛盾。]

\begin{remark}[7.7] 注意，如上讨论的剪切坐标给出了簇代数的最简单例子，其结构群是第 5 章详细研究的群 $\SL_2(\R)$。\end{remark}
